% STATUS: rewritten 2026-11-05
% STATUS: extended 2026-11-07 (Rakhimov and Al-Hossain 2012, Kurdachenko,
% Semko and Yashchuk 2024, Kurbanbaev 2025, Di Bartolo, Di Fatta and La Rosa
% 2025 and Kurdachenko and Semko 2025 read in full)
% Sources read in full: Loday and Pirashvili 1993 (Definition (1.1), p. 141),
% Ayupov and Omirov 1998 (p. 2), Biyogmam and Tcheka 2021 (Definition 3.3,
% Proposition 3.5, pp. 2192--2193), Kurdachenko--Pypka--Semko 2026 (pp. 1--2);
% Rakhimov and Al-Hossain 2012 (pp. 1--9); Kurdachenko, Semko and Yashchuk
% 2024 (pp. 819--822); Kurbanbaev 2025 (pp. 1525--1531); Di Bartolo, Di Fatta
% and La Rosa 2025 (pp. 1--9); Kurdachenko and Semko 2025 (pp. 110--114).
% Sources: LodayPirashvili1993, AyupovOmirov1998Engel,
% BiyogmamTcheka2021Outer, KurdachenkoPypkaSemko2026.
\section{Leibniz algebras}\label{sec:nonassociative:leibniz}

The previous sections asked for the Jordan or Malcev identities. % claim-ok: own summary of the chapter so far, links to Loday below
This section keeps only one Leibniz rule on the bracket. Loday introduced
the class in 1993, while he studied a noncommutative form of the
Chevalley--Eilenberg complex \cite[Introduction, p.~139]{LodayPirashvili1993}.

\begin{definition}[{\cite[Definition~(1.1),
p.~141]{LodayPirashvili1993}}]
\label{def:nonassociative:leibniz}
A \emph{Leibniz algebra} $\mathfrak{g}$ over a commutative ring $k$ is a
$k$-module with a bilinear bracket $[-,-]\colon\mathfrak{g}\times
\mathfrak{g}\to\mathfrak{g}$ satisfying the Leibniz identity
\[
[x,[y,z]]\;=\;[[x,y],z]-[[x,z],y]
\qquad\text{for all }x,y,z\in\mathfrak{g}.
\]
\end{definition}

This is the right-Leibniz convention, in which identity~(1.1) controls the
right factor. The left-Leibniz convention swaps the two factors
\cite[Definition~(1.1), p.~141]{LodayPirashvili1993}. A Leibniz algebra with
$[x,x]=0$ for every $x$ is a Lie algebra, since the Leibniz identity then
becomes the Jacobi identity \cite[Definition~(1.1),
p.~141]{LodayPirashvili1993}. Ayupov and Omirov note that the right
multiplications of a Leibniz algebra act like the inner derivations of a Lie
algebra \cite[p.~2]{AyupovOmirov1998Engel}.

\begin{proposition}[{\cite[p.~2]{AyupovOmirov1998Engel}}]
\label{prop:nonassociative:leibniz-right}
Let $L$ be a Leibniz algebra over a field $F$, and for $x\in L$ write
$R_x(z)=[z,x]$ for right multiplication by $x$. Then the set
$R(L)=\{R_x:x\in L\}$ with the commutator product is a Lie algebra, and
$R_xR_y-R_yR_x=R_{[y,x]}$ holds in $R(L)$.
\end{proposition}

Biyogmam and Tcheka measure the gap from the Lie case by one ideal
\cite[p.~2191]{BiyogmamTcheka2021Outer}.

\begin{definition}[{\cite[p.~2191]{BiyogmamTcheka2021Outer}}]
\label{def:nonassociative:leibniz-kernel}
For a Leibniz algebra $\mathfrak{g}$, Biyogmam and Tcheka write
$\mathrm{Leib}(\mathfrak{g})$ for the subspace spanned by all elements
$[x,x]$ with $x\in\mathfrak{g}$. It is a two-sided ideal of
$\mathfrak{g}$, called the \emph{Leibniz kernel}, and the quotient
$\mathfrak{g}_{\mathrm{Lie}}=\mathfrak{g}/\mathrm{Leib}(\mathfrak{g})$
is a Lie algebra.
\end{definition}

This kernel brings innerness and completeness to Leibniz algebras
\cite[Definition~2.4 and Definition~3.3,
p.~2192]{BiyogmamTcheka2021Outer}.

\begin{definition}[{\cite[Definition~2.4,
p.~2192]{BiyogmamTcheka2021Outer}}]
\label{def:nonassociative:leibniz-inner}
A derivation $d\colon\mathfrak{g}\to\mathfrak{g}$ of a Leibniz algebra
$\mathfrak{g}$ is an \emph{inner derivation} if
$d(\mathfrak{g})\subseteq R_x(\mathfrak{g})+
\mathrm{Leib}(\mathfrak{g})$ for some $x\in\mathfrak{g}$.
\end{definition}

A derivation that is not inner is an \emph{outer derivation}
\cite[Definition~3.3, p.~2192]{BiyogmamTcheka2021Outer}. Biyogmam and Tcheka
call $\mathfrak{g}$ \emph{complete} when
$Z(\mathfrak{g}/\mathrm{Leib}(\mathfrak{g}))=0$ and $\mathfrak{g}$ has no
outer derivation \cite[Definition~3.3, p.~2192]{BiyogmamTcheka2021Outer}.
Every semisimple Leibniz algebra in characteristic zero is complete. The
derivations themselves can be written down. On a simple complex Leibniz algebra
$L=G\dotplus I$, every derivation is $R_{a}+\alpha+\varphi$. Here $a\in G$,
the map $\alpha$ on the ideal respects the bracket, and $\varphi$ carries $G$
into that ideal \cite[Theorem~3.2]{RakhimovMasutovaOmirov2017}. Rakhimov,
Masutova and Omirov give the same description for a subclass of the semisimple
algebras \cite[Abstract, p.~1]{RakhimovMasutovaOmirov2017}.

\begin{proposition}[{\cite[Proposition~3.5,
p.~2193]{BiyogmamTcheka2021Outer}}]
\label{prop:nonassociative:leibniz-no-outer}
Let $\mathfrak{g}$ be a semisimple Leibniz algebra over a field of
characteristic zero. Then $\mathfrak{g}$ admits no outer derivations.
More precisely, each derivation $d$ of $\mathfrak{g}$ can be written as
$d=R_z+a$ for some $z\in\mathfrak{g}$ and some derivation
$a\colon\mathfrak{g}\to \mathrm{Leib}(\mathfrak{g})$ that vanishes on
the Levi subalgebra of $\mathfrak{g}$.
\end{proposition}

At dimension three the description is complete. Kurdachenko, Pypka and Semko
take the non-Lie left Leibniz algebras of dimension three over any field
\cite[p.~1]{KurdachenkoPypkaSemko2026}. Earlier work split them into sixteen
types, and the authors list the derivations of each type by direct calculation
\cite[p.~1]{KurdachenkoPypkaSemko2026}.

\begin{theorem}[{\cite[Theorem~1.1, p.~2]{KurdachenkoPypkaSemko2026}}]\label{thm:nonassociative:leibniz-dim3}
Let $L$ be a three-dimensional non-Lie left Leibniz algebra over an
arbitrary field $F$. Then its derivation algebra $\operatorname{Der}(L)$ is
the subspace of $M_{3}(F)$ given by the formulas proved or quoted in the
cited paper and collected in its Table~2. The table covers every parameter
family and every case of characteristic two in the sixteen-type list.
\end{theorem}

The table shows that one derivation algebra grows in characteristic two. In
that characteristic several types gain derivations
\cite[p.~1]{KurdachenkoPypkaSemko2026}. So the dimension can change along a
family. The authors also record how the resulting Lie algebras split into
ideals and subalgebras \cite[p.~1]{KurdachenkoPypkaSemko2026}.

Nilpotent Leibniz algebras of class three also have an explicit answer at
dimension three. Kurdachenko, Semko and Yashchuk take the algebras of dimension
three whose class is three, and those whose centre has dimension two
\cite[p.~819]{KurdachenkoSemkoYashchuk2024}. Take the first algebra. Its
derivation algebra is a semidirect sum. The ideal is abelian. The complement is
generated by one derivation, which scales the three basis vectors by $1$, $2$
and $3$ \cite[Theorem~3.1, pp.~821--822]{KurdachenkoSemkoYashchuk2024}.

The derivation algebra also has a matrix shape. It is isomorphic to the Lie
algebra of matrices
\[
\begin{pmatrix}\alpha_1&0&0\\ \alpha_2&2\alpha_1&0\\ \alpha_3&\alpha_2&3\alpha_1\end{pmatrix},
\qquad \alpha_1,\alpha_2,\alpha_3\in F,
\]
so it is three-dimensional
\cite[Theorem~3.1, pp.~821--822]{KurdachenkoSemkoYashchuk2024}. The second algebra
of the paper has a derivation algebra of another shape
\cite[Theorem~3.2, p.~825]{KurdachenkoSemkoYashchuk2024}.

The filiform families have a dimension count that depends on the structure
constants. Rakhimov and Al-Hossain take three naturally graded complex filiform
algebras. They also take two further families that come from them
\cite[Theorem~1, p.~3]{RakhimovAlHossain2012}. They compute the dimension of
the derivation algebra in each family
\cite[Theorems~3--5, pp.~4, 8 and~12]{RakhimovAlHossain2012}.

\begin{theorem}[{\cite[Theorems~3--5, pp.~4, 8 and~12]{RakhimovAlHossain2012}}]\label{thm:nonassociative:filiform-dim}
Let $n\geq2$. The derivation algebras of $NGF_1$, $NGF_2$ and $NGF_3$ have
dimensions $n+1$, $n+2$ and $2n-1$. For $L$ in $FL_b^n$ the dimension lies
between $n-1$ and $n+1$, and for $L$ in $SL_b^n$ it lies between $n-1$ and
$n+2$.
\end{theorem}

So the dimension is not fixed by the filiform class. The structure constants
of each family decide it \cite[pp.~8--12]{RakhimovAlHossain2012}.

Almost inner derivations ask for less than innerness. A derivation $D$ of a
Leibniz algebra $L$ is almost inner when every $x\in L$ has some $a_x\in L$
with $D(x)=[x,a_x]$ \cite[Definition~2.4, p.~1527]{Kurbanbaev2025}. Every inner
derivation is almost inner. The question is whether the two spaces agree
\cite[p.~1527]{Kurbanbaev2025}. On a simple Leibniz algebra the answer depends
on the shape of the algebra.

\begin{theorem}[{\cite[Theorems~4.1--4.2, pp.~1530--1531]{Kurbanbaev2025}}]\label{thm:nonassociative:almost-inner}
Let $L=\mathfrak{g}\ltimes I$ be a simple Leibniz algebra over $\mathbb{C}$.
\begin{enumerate}
\item If $\mathfrak{g}$ and $I$ are isomorphic as $\mathfrak{g}$-modules, then
every almost inner derivation of $L$ is inner.
\item If they are not isomorphic, then every almost inner derivation of $L$ is
inner when $\mathfrak{g}\cong\mathfrak{sl}_2$, and there is an almost inner
derivation that is not inner otherwise.
\end{enumerate}
\end{theorem}

The proof tells the two cases apart by one strongly regular element of a
Cartan subalgebra \cite[p.~1530]{Kurbanbaev2025}. So the exceptional case
belongs to the smallest simple Lie algebra.

Complete Leibniz algebras settle the bilinear form of the same question. The
literature has two definitions of completeness, and they differ
\cite[Definitions~2.6--2.7, p.~6]{DiBartoloDiFattaLaRosa2025Leibniz}. Under the
first, $L$ is complete when $Z(L/\mathrm{Leib}(L))=0$ and every derivation is
inner up to an element of $\mathrm{Leib}(L)$
\cite[Definition~2.6, p.~6]{DiBartoloDiFattaLaRosa2025Leibniz}. Under the
second, $L$ is complete when $Z(L)=0$ and every derivation of $L$ is inner
\cite[Definition~2.7, p.~6]{DiBartoloDiFattaLaRosa2025Leibniz}.

\begin{theorem}[{\cite[Theorems~3.1--3.2, pp.~7--8]{DiBartoloDiFattaLaRosa2025Leibniz}}]\label{thm:nonassociative:complete-biderivation}
Let $L$ be a complete Leibniz algebra. Then a bilinear map $f\colon L\times
L\to L$ is a biderivation of $L$ if and only if there are linear maps
$\varphi,\psi\colon L\to L$ with $f(x,y)=[\varphi(x),y]=[\psi(y),x]$ for all
$x,y\in L$, in the sense of Definition~2.7. In the sense of Definition~2.6 the
same holds up to a right and a left biderivation with values in
$\mathrm{Leib}(L)$.
\end{theorem}

The two versions differ exactly in the Leibniz kernel
\cite[pp.~7--8]{DiBartoloDiFattaLaRosa2025Leibniz}. One example is complete in
one sense and not in the other \cite[Example~2.7.1, p.~6]{DiBartoloDiFattaLaRosa2025Leibniz}.
So the two statements are not interchangeable.

The non-nilpotent side has a calculation of its own. Kurdachenko, Semko and
Subbotin take some non-nilpotent Leibniz algebras of low dimension and list the
derivations of each \cite[Theorems~1--3, pp.~120--126]{KurdachenkoSemkoSubbotin2025}.
A second paper does the same for further families
\cite[Theorems~1--2, pp.~209--211]{KurdachenkoSemkoSubbotin2025b}.

The next section separates the even and odd parts of an algebra: % claim-ok: own signpost to the superalgebra section
it treats derivations that respect a grading by
$\mathbb{Z}/2\mathbb{Z}$.
